\documentclass[11pt,a4paper]{amsart}
% Adapted from an OpenAI Math manuscript. See PROVENANCE.md.

\usepackage{cgeometry}
\crefname{maintheorem}{Theorem}{Theorems}
\Crefname{maintheorem}{Theorem}{Theorems}

\addbibresource{ComplexGeometry.bib}
\defbibheading{bibliography}[\refname]{\section*{#1}}
\setlength{\bibitemsep}{0pt}
\renewcommand*{\bibfont}{\footnotesize}

\allowdisplaybreaks[2]
\emergencystretch=1em

%%%%%%%%%%%%%%%%%%%%%%%%%%%%%%%%%%%%%%%%%%%%%%%%%%%%%%%%%%%%%%%%%%%%%%%%%%%%%
%% Notation
%%%%%%%%%%%%%%%%%%%%%%%%%%%%%%%%%%%%%%%%%%%%%%%%%%%%%%%%%%%%%%%%%%%%%%%%%%%%%



\DeclareMathOperator{\supp}{supp}

\title{The logarithmic Brunn--Minkowski conjecture}
\author{}
\date{}

\begin{document}

\begin{abstract}
We prove the logarithmic Brunn--Minkowski inequality for origin-symmetric convex bodies in every dimension. It implies the symmetric $L_p$ Brunn--Minkowski inequality for $0<p<1$ and, by the transfer theorem of Saroglou, the scalar-dilation $(B)$-conjecture for even log-concave Radon measures.
\end{abstract}

\maketitle


% Adapted from an OpenAI Math manuscript. See PROVENANCE.md.
\section{Introduction}
\label{sec:introduction}

A convex body in $\mathbb{R}^n$ is a compact convex set with nonempty interior. If $K=-K$, then $0$ lies in its interior. Its support function is
\[
 h_K(u)=\max_{x\in K}\langle x,u\rangle,\qquad u\in S^{n-1}.
\]
For an origin-symmetric body this function is strictly positive. For a positive continuous function $f$ on the unit sphere, its \emph{Wulff body} is
\[
 \mathcal{W}[f]=\bigcap_{u\in S^{n-1}}
 \{x\in\mathbb{R}^n:\langle x,u\rangle\le f(u)\}.
\]
It is a convex body: it contains the ball of radius $\min f>0$ and is contained in the ball of radius $\max f$. Write $|K|$ for $n$-dimensional Lebesgue measure. Our main result is the following.

\begin{maintheorem}
\label{thm:main}
Let $n\ge1$ and let $K,L\subset\mathbb{R}^n$ be convex bodies satisfying $K=-K$ and $L=-L$. For every $0\le\lambda\le1$,
\begin{equation}\label{eq:main}
 \left|\mathcal{W}[h_K^{1-\lambda}h_L^\lambda]\right|
 \ge |K|^{1-\lambda}|L|^\lambda.
\end{equation}
\end{maintheorem}

This exposition follows the proof in \cite{OpenAIMathSource}.

The conjecture was formulated by B\"or\"oczky, Lutwak, Yang, and Zhang, who proved it in dimension two \cite[Problem~1.1 and Theorem~1.3]{BLZY2012}. Saroglou proved the inequality for pairs unconditional in a common basis \cite[Theorem~1.2]{Saroglou2015}. The proof below uses evenness through a density argument for the centered test functions.

\subsection{A variance inequality for slab widths}

For a finite spanning family of unit vectors $p_i$ and positive slab widths $h_i$, approximate the intersection of the strips $|\langle p_i,x\rangle|\leq h_i$ by the weight
\[
 \exp\left(-\frac{\varepsilon}{2}|x|^2
       -\sum_i\left(\frac{\langle p_i,x\rangle}{h_i}\right)^q\right),
 \qquad \varepsilon>0,\quad q\geq2\text{ even}.
\]
The successive limits $q\to\infty$ and $\varepsilon\downarrow0$ recover the volume of the strip intersection. Approximating the support bounds by finite sets of directions then reduces \eqref{eq:main} to log-concavity of the smooth integral as the widths vary geometrically (\cref{sec:reduction}).

The second derivative of the log integral is a variance minus the mean of a second derivative. In moment coordinates $x=\nabla\varphi(z)$, the probability measure for this variance is the pushforward of $\ee^{-\varphi(z)}\,\rd z$. The Hessian of $\varphi$ has a global upper bound and is positive definite everywhere (\cref{sec:moment}). A divergence operator built from this Hessian has dense image in the centered even test space. The affine obstructions are removed by centering and evenness (\cref{sec:identities,sec:density}). A sum-of-squares tensor identity then gives the summandwise variance estimate needed for the concavity (\cref{sec:tensor,sec:variance}).

\subsection{Consequences}

The logarithmic endpoint gives the additive $L_p$ inequality by normalizing the two volumes and reweighting the interpolation parameter. Saroglou's transfer theorem then gives the logarithmic inequality for even log-concave densities \cite[Theorem~3.1]{Saroglou2016}. \Cref{sec:consequences} proves the following corollary and the scalar-dilation $(B)$-conjecture for even log-concave Radon measures, including measures supported on proper subspaces.

\begin{corollary}
\label{cor:lp}
Let $n\ge1$, let $K,L\subset\mathbb{R}^n$ be full-dimensional origin-symmetric convex bodies, and let $0<p<1$. For every $0\le\lambda\le1$,
\begin{equation}\label{eq:lp}
 \left|\mathcal{W}\!\left[
  \bigl((1-\lambda)h_K^p+\lambda h_L^p\bigr)^{1/p}
 \right]\right|^{p/n}
 \ge (1-\lambda)|K|^{p/n}+\lambda|L|^{p/n}.
\end{equation}
\end{corollary}

% Adapted from an OpenAI Math manuscript. See PROVENANCE.md.
\section{From a variance bound to volume}
\label{sec:reduction}

We first deduce \cref{thm:main} from the variance bound below. For a probability measure $\mu$, write $\mathbb{E}_\mu f=\int f\,\rd \mu$ and $\Var_\mu(f)=\mathbb{E}_\mu(f-\mathbb{E}_\mu f)^2$.

\begin{proposition}\label{prop:variance}
Let $n\ge2$, $N\ge1$, $q\ge2$ be an even integer, and $\varepsilon>0$. Let $p_1,\ldots,p_N\in\mathbb{R}^n\setminus\{0\}$ and $h_1,\ldots,h_N>0$. Set
\begin{equation}\label{eq:potential}
 \psi_i(x)=\left(\frac{\langle p_i,x\rangle}{h_i}\right)^q,
 \qquad V(x)=\frac{\varepsilon}{2}|x|^2+\sum_{i=1}^N\psi_i(x),
 \qquad d\mu=Z^{-1}\ee^{-V(x)}\,\rd x,
\end{equation}
where $Z=\int_{\mathbb{R}^n}\ee^{-V(x)}\,\rd x$. For all real $b_1,\ldots,b_N$,
\begin{equation}\label{eq:variance}
 \Var_\mu\left(\sum_{i=1}^N b_i\psi_i\right)
 \le\sum_{i=1}^N b_i^2\mathbb{E}_\mu\psi_i.
\end{equation}
\end{proposition}

Here the Gaussian term makes $Z$ finite, even if the vectors $p_i$ do not span $\mathbb{R}^n$. All polynomial moments are finite. The proof of \cref{prop:variance} occupies Sections~\ref{sec:moment}--\ref{sec:variance}.


The proof will establish the dual estimate
\[
 \sum_{i=1}^N\frac{(\mathbb{E}_\mu u\psi_i)^2}{\mathbb{E}_\mu\psi_i}
 \le \mathbb{E}_\mu u^2
 \qquad\text{for every centered even }u\in L^2(\mu).
\]
Here centered means $\mathbb{E}_\mu u=0$, and the denominators are positive because each $p_i\ne0$ and $\mu$ has a positive density. Applied to $u=\sum_i b_i\psi_i-\mathbb{E}_\mu\sum_i b_i\psi_i$, this gives \eqref{eq:variance} by Cauchy--Schwarz.  We first prove the dual estimate on a class of smooth compactly supported tests, then use density to reach every such $u$.  This explains the roles of the density and tensor lemmas in the analytic proof.

\begin{proof}[Proof of \cref{thm:main} from \cref{prop:variance}]
Suppose first that $n\ge2$. Fix nonzero spanning vectors $p_1,\ldots,p_N$ and positive endpoint widths $h_i^0,h_i^1$. For $t\in[0,1]$, define
\[
 h_i(t)=(h_i^0)^{1-t}(h_i^1)^t,\qquad
 P(t)=\{x:|\langle p_i,x\rangle|\le h_i(t)\text{ for all }i\}.
\]
The spanning condition makes each $P(t)$ bounded; positivity of the widths gives nonempty interior. Use these widths in \eqref{eq:potential} and denote the resulting integral by $Z_{q,\varepsilon}(t)$. For fixed $q,\varepsilon$, its first two derivatives may be taken under the integral: on a compact interval of $t$, the differentiated integrands are bounded by a polynomial times $\ee^{-\varepsilon|x|^2/2}$. With $b_i=-q\log(h_i^1/h_i^0)$ one has
\[
 \dot V=\sum_i b_i\psi_i,\qquad
 \ddot V=\sum_i b_i^2\psi_i,
\]
and hence
\[
 \frac{d^2}{dt^2}\log Z_{q,\varepsilon}(t)
 =\Var_{\mu_t}(\dot V)-\mathbb{E}_{\mu_t}\ddot V\le0
\]
by \cref{prop:variance}. Therefore
\begin{equation}\label{eq:partition-concavity}
 Z_{q,\varepsilon}(\lambda)\ge
 Z_{q,\varepsilon}(0)^{1-\lambda}Z_{q,\varepsilon}(1)^\lambda.
\end{equation}

Keep $\varepsilon>0$ fixed and let $q$ tend to infinity through even integers. For each fixed $t$, the integrands converge almost everywhere to $\ee^{-\varepsilon|x|^2/2}\mathbf1_{P(t)}$. The exceptional set lies in the finitely many hyperplanes $\langle p_i,x\rangle=\pm h_i(t)$ and has measure zero. The common integrable bound $\ee^{-\varepsilon|x|^2/2}$ gives
\[
 Z_{q,\varepsilon}(t)\longrightarrow
 \int_{P(t)}\ee^{-\varepsilon|x|^2/2}\,\rd x.
\]
Apply this at $t=0,\lambda,1$ in \eqref{eq:partition-concavity}, then let $\varepsilon\downarrow0$. Monotone convergence gives
\begin{equation}\label{eq:slabs}
 |P(\lambda)|\ge |P(0)|^{1-\lambda}|P(1)|^\lambda.
\end{equation}

% Adapted from an OpenAI Math manuscript. See PROVENANCE.md.
\begin{figure}[tb]
  \centering
  \begin{tikzpicture}[x=1cm,y=1cm,line join=round]
    \definecolor{slabblue}{RGB}{34,100,139}
    \definecolor{slabfill}{RGB}{222,237,245}
    \begin{scope}[xshift=-2.8cm]
      \filldraw[fill=slabfill,draw=slabblue,line width=0.8pt]
        (-1.6,-1) rectangle (1.6,1);
      \filldraw[fill=white,draw=black,line width=0.6pt]
        (0,0) ellipse [x radius=1.6,y radius=1];
      \node at (0,0) {$E$};
      \node[anchor=north] at (0,-1.22) {$P_2$};
      \node[anchor=north,font=\small] at (0,-1.62)
        {Two slab directions};
    \end{scope}
    \begin{scope}[xshift=2.8cm]
      \foreach \k in {0,...,7} {
        \coordinate (p\k) at
          ({1.6*cos(22.5+45*\k)/cos(22.5)},
           {sin(22.5+45*\k)/cos(22.5)});
      }
      \filldraw[fill=slabfill,draw=slabblue,line width=0.8pt]
        (p0) \foreach \k in {1,...,7} {-- (p\k)} -- cycle;
      \filldraw[fill=white,draw=black,line width=0.6pt]
        (0,0) ellipse [x radius=1.6,y radius=1];
      \node at (0,0) {$E$};
      \node[anchor=north] at (0,-1.22) {$P_4$};
      \node[anchor=north,font=\small] at (0,-1.62)
        {Four slab directions};
    \end{scope}
  \end{tikzpicture}
  \caption{Finite outer approximations of an ellipse $E$. Adding slab
  directions gives $P_2\supset P_4\supset E$; the subscripts here count
  slab directions. Using nested finite spanning direction
  sets with dense union, the same construction recovers $\mathcal{W}[f]$ from any
  positive continuous even function $f$ on the sphere; $f$ need not be a
  support function.}
  \label{fig:slab-exhaustion}
\end{figure}


\Cref{fig:slab-exhaustion} illustrates how additional slab directions refine an outer approximation.  To pass to all directions, let $K,L$ be as in \cref{thm:main}. Choose nested finite sets $D_m\subset S^{n-1}$ that contain the coordinate vectors and whose union is dense in the sphere. Define
\[
 P_m(t)=\bigcap_{p\in D_m}
 \{x:|\langle p,x\rangle|\le h_K(p)^{1-t}h_L(p)^t\}.
\]
The bodies $P_m(t)$ decrease with $m$ and are all contained in the bounded body $P_1(t)$. By continuity of the support functions and their evenness,
\[
 \bigcap_m P_m(t)=\mathcal{W}[h_K^{1-t}h_L^t].
\]
Indeed any point satisfying the dense set of inequalities satisfies all of them by continuity in the direction. Conversely, the Wulff inequalities imply the absolute-value constraints by evenness. At $t=0,1$ these intersections are $K,L$, respectively, by the supporting-half-space characterization of a convex body. Continuity of measure from above, applied at $0,\lambda,1$ in \eqref{eq:slabs}, proves \eqref{eq:main}.

Finally, in dimension one write $K=[-a,a]$ and $L=[-b,b]$ with $a,b>0$. Their Wulff interpolation is the interval with radius $a^{1-\lambda}b^\lambda$, whose length is $(2a)^{1-\lambda}(2b)^\lambda$. Thus \eqref{eq:main} is equality. The endpoints $\lambda=0,1$ are equality in every dimension.
\end{proof}

% Adapted from an OpenAI Math manuscript. See PROVENANCE.md.
\section{Moment coordinates on the whole space}\label{sec:moment}

The variance argument uses a change of variables whose Jacobian is a positive Hessian.  The following lemma supplies this change of variables and the uniform upper bound needed for the later cutoff arguments.

\begin{lemma}\label{lem:moment}
Let $n\ge2$, $\varepsilon>0$, and let $V\in C^\infty(\mathbb{R}^n)$ be even with $D^2V\ge\varepsilon I$.  Set
\[
 Z=\int_{\mathbb{R}^n}\ee^{-V(x)}\,\rd x,
 \qquad d\mu(x)=Z^{-1}\ee^{-V(x)}\,\rd x.
\]
There is an even $\varphi\in C^\infty(\mathbb{R}^n)$ such that $d\nu(z)=\ee^{-\varphi(z)}\,\rd z$ is a probability measure, $\nabla\varphi$ is a diffeomorphism of $\mathbb{R}^n$ onto $\mathbb{R}^n$, and
\begin{equation}\label{eq:moment}
 (\nabla\varphi)_\#\nu=\mu,
 \qquad 0<D^2\varphi\le\frac4\varepsilon I,
 \qquad
 \log\det D^2\varphi
   =V(\nabla\varphi)+\log Z-\varphi.
\end{equation}
Moreover, $\varphi(z)\ge c|z|-b$ for some $c>0$ and $b<\infty$. The lower Hessian bound in \eqref{eq:moment} means positive definiteness at every point; it need not be uniform on $\mathbb{R}^n$.
\end{lemma}

\begin{proof}
Evenness and uniform convexity give $V(x)\ge V(0)+\varepsilon|x|^2/2$, so $0<Z<\infty$. For $R\ge2$, write $B_R=\{x:|x|<R\}$ and set
\[
 Z_R=\int_{B_R}\ee^{-V(x)}\,\rd x,
 \qquad d\mu_R=Z_R^{-1}\mathbf1_{B_R}\ee^{-V(x)}\,\rd x.
\]
We first use two compact-support results.  The moment theorem of Berman--Berndtsson \cite[Theorem~1.1]{BermanBerndtsson2013}, applied to the body $\overline B_R$ and the positive smooth density $\ee^{-V}/Z_R$, gives a smooth convex potential $\varphi_R$, unique up to translation, for which $\nabla\varphi_R:\mathbb{R}^n\to B_R$ is a diffeomorphism and
\begin{equation}\label{eq:moment-truncated}
 \det D^2\varphi_R
   =Z_R\exp\{V(\nabla\varphi_R)-\varphi_R\}.
\end{equation}
The centering hypothesis holds by evenness.  Integrating the equation through this diffeomorphism shows that $d\nu_R=\ee^{-\varphi_R}\,\rd z$ is a probability and $(\nabla\varphi_R)_\#\nu_R=\mu_R$. Klartag's estimate \cite[Proposition~2 and Remark~1]{Klartag2014Moment} gives
\[
 0<D^2\varphi_R\le CI,\qquad C=4/\varepsilon.
\]
Indeed, the ball has smooth positively curved boundary, and $\rho_R=V+\log Z_R$ is convex and smooth, with itself and all its derivatives bounded on $B_R$, and $D^2\rho_R\ge\varepsilon I$. The constant $C$ is independent of $R$.

Translate the unique preimage of $0$ to $0$.  It is the unique minimum of $\varphi_R$.  Reflection gives a second potential for the same target with its minimum at $0$; uniqueness up to translation therefore makes $\varphi_R$ even.  For each fixed $R$, take the preimages of the slopes $\pm(R/2)e_i$.  Their supporting planes give
\[
 \varphi_R(z)\ge\frac{R}{2\sqrt n}|z|-b_R
\]
with a finite, possibly $R$-dependent constant $b_R$. Thus $\nu_R$ has an exponential tail.  In particular, since $|\nabla\varphi_R|\le R$, integration of $\operatorname{div}(z \ee^{-\varphi_R})$ with expanding compact cutoffs is justified and yields
\begin{equation}\label{eq:moment-source-identity}
 \mathbb{E}_{\nu_R}\,z\cdot\nabla\varphi_R(z)=n.
\end{equation}

We now obtain bounds independent of $R$.  Put $m_R=\varphi_R(0)$ and $g_R=\varphi_R-m_R\ge0$. The Hessian bound implies $g_R(y)\le C|y|^2/2$ and $|\nabla\varphi_R(y)|\le C|y|$. For arbitrary $y,z$, set $x=\nabla\varphi_R(z)$.  Supporting planes at $z$ and $-z$ give the single pointwise inequality
\[
 g_R(y)\ge g_R(z)+|x\cdot y|-x\cdot z.
\]
Integrate with respect to $\nu_R$, discard the nonnegative $\mathbb{E}_{\nu_R}g_R$, and use \eqref{eq:moment-source-identity}:
\[
 g_R(y)\ge\int|x\cdot y|\,\rd \mu_R(x)-n.
\]
On $B_1$ the density of $\mu_R$ is at least $d_0=Z^{-1}\exp(-\max_{\overline B_1}V)>0$. Rotational symmetry of the ball therefore bounds the integral below by $c|y|$, where $c=d_0\int_{B_1}|x_1|\,\rd x>0$.  Hence
\begin{equation}\label{eq:moment-coercivity}
 c|y|-n\le g_R(y)\le\frac C2|y|^2,
 \qquad
 \left(\frac{2\pi}{C}\right)^{n/2}
 \le \ee^{m_R}=\int \ee^{-g_R(y)}\,\rd y
 \le \ee^n\int \ee^{-c|y|}\,\rd y.
\end{equation}
This bounds $m_R$ on both sides and gives one integrable exponential majorant for all the densities $\ee^{-\varphi_R}$.

It remains to obtain a smooth limit and to verify that its gradient reaches every point.  On each fixed ball, the preceding estimates bound $\varphi_R$ and $\nabla\varphi_R$ uniformly.  Since $Z_2\le Z_R\le Z$, equation \eqref{eq:moment-truncated} bounds $\det D^2\varphi_R$ below by a positive constant there. Together with $D^2\varphi_R\le CI$, this gives a uniform local bound $\delta I\le D^2\varphi_R\le CI$: if the determinant is at least $d>0$, every eigenvalue is at least $d/C^{n-1}$. The right-hand sides
\[
 f_R=V(\nabla\varphi_R)+\log Z_R-\varphi_R
\]
are uniformly Lipschitz on that ball. Indeed, $\nabla\varphi_R$ takes values in a fixed compact set there, and
\[
 \nabla f_R=D^2\varphi_R\nabla V(\nabla\varphi_R)-\nabla\varphi_R
\]
is uniformly bounded.

To apply an interior uniformly elliptic estimate, choose a smooth concave scalar function $s$ equal to $\log$ on a slightly larger interval than $[\delta,C]$, with $s'$ bounded above and bounded below by positive constants on $\mathbb{R}$.  Such an extension is obtained by smoothly continuing the decreasing derivative $1/t$ to positive constant tails.  For symmetric matrices $Q$, the spectral function $\mathcal{F}(Q)=\tr s(Q)$ is smooth, concave and uniformly elliptic: its derivative has eigenvalues $s'(\lambda_i(Q))$, and its second variation has coefficients $s''(\lambda_i)$ and the nonpositive divided differences of $s'$.  It agrees with $\log\det$ on the Hessians under consideration.  The nonhomogeneous Evans--Krylov estimate \cite[Corollary~2.3]{Wang2012} applied to $\mathcal{F}(D^2\varphi_R)=f_R$ gives uniform interior $C^{2,\alpha}$ bounds for some $\alpha>0$.

Each $\varphi_R$ is already smooth.  Its derivative $w_{R,k}=\partial_k\varphi_R$ satisfies
\[
 (D^2\varphi_R)^{-1}:D^2w_{R,k}
   =\nabla V(\nabla\varphi_R)\cdot\nabla w_{R,k}-w_{R,k},
\]
where $Q:N=\tr(Q^TN)$.  The principal coefficients and the right-hand side are uniformly $C^\alpha$ on smaller balls. Interior Schauder estimates \cite[Theorem~2.20 and Corollary~2.21]{FernandezRealRosOton2022} give uniform $C^{3,\alpha}$ bounds for $\varphi_R$, and iteration gives bounds for every derivative on compact sets. A diagonal subsequence, still indexed by $R$, converges locally smoothly to an even $\varphi$ with positive definite Hessian, the same upper cap, and the logarithmic equation in \eqref{eq:moment}. The bound $\varphi_R(z)\ge c|z|-b$, with $b$ independent of $R$, also passes to $\varphi$.

The common exponential majorant from \eqref{eq:moment-coercivity} proves $\int \ee^{-\varphi}=1$ by dominated convergence.  For each bounded continuous $h$, it also gives
\[
 \int h(\nabla\varphi)\ee^{-\varphi}
 =\lim_{R\to\infty}\int h(\nabla\varphi_R)\ee^{-\varphi_R}
 =\lim_{R\to\infty}\int h\,\rd \mu_R
 =\int h\,\rd \mu.
\]
Thus $(\nabla\varphi)_\#\nu=\mu$. Positive definite Hessian makes the gradient locally invertible and injective, since for $z\ne w$,
\[
 (\nabla\varphi(z)-\nabla\varphi(w))\cdot(z-w)
 =\int_0^1 D^2\varphi(w+t(z-w))[z-w,z-w]\,\rd t>0.
\]
Its image is dense: an open ball missed by the image would have zero pushforward mass, contrary to the everywhere positive density of $\mu$.  Fix $x_0\in\mathbb{R}^n$.  Density permits finitely many image points $y_j=\nabla\varphi(z_j)$ whose convex hull contains a ball $B_r(x_0)$ with $r>0$; one may take small perturbations of the points $x_0\pm e_i$.  Their supporting planes imply
\[
 \varphi(z)-x_0\cdot z
 \ge\max_j\{(y_j-x_0)\cdot z+\varphi(z_j)-y_j\cdot z_j\}
 \ge r|z|-b_{x_0}.
\]
This function attains a minimum, at which $\nabla\varphi=x_0$.  The gradient is therefore onto, and its local smooth inverses combine to a global smooth inverse.
\end{proof}

We will write $x=\nabla\varphi(z)$ and $\tau=D_z^2\varphi$, evaluating $\tau$ at the corresponding point when using $x$ coordinates.  Thus $\tau(x)=D_z^2\varphi((\nabla\varphi)^{-1}(x))$ is smooth, even, positive definite and bounded above by $(4/\varepsilon)I$. The gradient is odd, as is its inverse.  Finally, a compactly supported smooth function of $x$ is compactly supported after composition with $\nabla\varphi$: its support lies in the image of a compact set under the continuous inverse map.  Consequently the compact integrations by parts below are valid in either coordinate system.

% Adapted from an OpenAI Math manuscript. See PROVENANCE.md.
\section{Adjoints in moment coordinates}\label{sec:identities}

We retain the probability measures
\[
 d\mu(x)=Z^{-1}\ee^{-V(x)}\,\rd x,
 \qquad d\nu(z)=\ee^{-\varphi(z)}\,\rd z,
 \qquad x=\nabla\varphi(z),
\]
and the smooth diffeomorphism constructed above.  Put
\[
 \tau=D_z^2\varphi,\qquad M=\tau^{-1}.
\]
In target coordinates, these symbols denote the same matrix fields evaluated at $z=(\nabla\varphi)^{-1}(x)$. Expectation under either measure, with functions identified by this map, is denoted by $\mathbb{E}$. Throughout the coordinate calculations, repeated indices are summed from $1$ to $n$, and $Q:L=\tr(Q^TL)$ denotes matrix contraction.

We first construct the compact test functions used in the variance argument and establish an identity that will also identify the obstruction to their density.  Write
\[
 \partial_k=\partial_{z_k},\qquad
 \partial_k^*=x_k-\partial_k.
\]
The latter is the adjoint of $\partial_k$ in $L^2(\nu)$ because $\partial_k\varphi=x_k$.  In target coordinates, write $\delta_\mu$ for negative weighted divergence:
\[
 \delta_\mu W
 =\sum_j\bigl(V_{x_j}W_j-\partial_{x_j}W_j\bigr).
\]
On a matrix, $\delta_\mu$ acts on each row, producing a vector. For a smooth row field $P_{ak}$ and a compact test $g$, the chain rule and adjunction give
\[
 \mathbb{E}g\,\partial_k^*P_{ak}
 =\mathbb{E}(\partial_k g)P_{ak}
 =\mathbb{E}g_{x_i}\tau_{ik}P_{ak}.
\]
Consequently
\begin{equation}\label{eq:row-adjoint}
 \partial_k^*P_{ak}=(\delta_\mu(P\tau))_a.
\end{equation}
Taking $P$ to be the identity matrix also gives $\delta_\mu\tau=x$. This is Fathi's moment-map construction of a Stein kernel \cite[Theorem~2.3 of the cited arXiv version]{Fathi2019}: the matrix field $\tau$ expresses integration by parts against $x$ through $\mathbb{E}x_j g=\mathbb{E}\tau_{jk}g_{x_k}$.

Thus the differential operator
\begin{equation}\label{eq:operator-A}
 Af=\delta_\mu(\tau\nabla_x f)
   =\partial_k^*f_{x_k}
   =x\cdot\nabla_x f-\tau:D_x^2f
\end{equation}
is formally symmetric, with
\begin{equation}\label{eq:A-form}
 \mathbb{E}gAf=\mathbb{E}\nabla_xg\cdot\tau\nabla_xf
\end{equation}
for compact smooth tests.  These formulas involve no operator-domain assertion. This is the negative of the moment-Hessian diffusion generator studied by Klartag \cite[equation~(68)]{Klartag2014Moment}; for our sign convention, $Ax_j=x_j$.

For $f\in C_c^\infty(\mathbb{R}_x^n)$, define
\begin{equation}\label{eq:test-fields}
 h=D_x^2f,\qquad P=\tau h,\qquad B=\tau h\tau,
 \qquad W_a=\partial_k^*P_{ak},\qquad u=\delta_\mu W.
\end{equation}
All these fields have compact support in both coordinates: the inverse image of a compact target set is its compact image under the continuous inverse diffeomorphism.  By \eqref{eq:row-adjoint}, $W=\delta_\mu B$.  For every smooth function $F$ in target coordinates, two adjunctions give
\begin{equation}\label{eq:test-pairing}
 \mathbb{E}uF=\mathbb{E}W\cdot\nabla_xF=\mathbb{E}B:D_x^2F.
\end{equation}
The fields are compactly supported, so $F$ need not be. Thus the same pairing can be estimated through either the first or the second derivatives of $F$. To identify the range of these tests, we also use the gradient identity
\begin{equation}\label{eq:W-gradient}
 W=\nabla_z\bigl[((A-1)f)(x(z))\bigr]
   =\tau\nabla_x(A-1)f,
 \qquad u=A(A-1)f.
\end{equation}
To check it directly, put $C_{ijk}=\partial_{z_i}\partial_{z_j}\partial_{z_k}\varphi$. This tensor is fully symmetric, and expansion gives
\[
 W_a=x_k\tau_{ai}f_{x_ix_k}
       -C_{aik}f_{x_ix_k}
       -\tau_{ai}\tau_{jk}f_{x_ix_kx_j}.
\]
Differentiating $(A-1)f$ in $z_a$ gives exactly these terms: the first-gradient terms from $x\cdot\nabla_xf$ and $-f$ cancel. The formula for $u$ follows from \eqref{eq:operator-A}. In particular, $D_zW$ is symmetric.

The following identity is the moment-coordinate specialization of the transport-Hessian Bochner formula of Kolesnikov, Livshyts and Rotem \cite[Proposition~4.8]{KLR2026}. We give its direct derivation in the present notation. Two compact integrations by parts give
\begin{equation}\label{eq:hessian}
 \begin{aligned}
 \mathbb{E}(Af)(A-1)f
 &=\mathbb{E}\nabla_x f\cdot\tau\nabla_x(A-1)f
   =\mathbb{E}\nabla_x f\cdot W\\
 &=\mathbb{E}h:B
   =\mathbb{E}\bigl\|\tau^{1/2}D_x^2f\,\tau^{1/2}\bigr\|_{\mathrm{HS}}^2.
 \end{aligned}
\end{equation}
Here $\|Q\|_{\mathrm{HS}}^2=\tr(Q^TQ)$. The differential expression $A$ satisfies $A1=0$ and $Ax_j=x_j$, so $A(A-1)$ annihilates every affine function. \Cref{sec:density} uses \eqref{eq:hessian} to prove that these are the only $L^2$ obstructions to density. Centering removes constants, and evenness removes linear functions.

% Adapted from an OpenAI Math manuscript. See PROVENANCE.md.
\section{Density of the even tests}\label{sec:density}

The Hessian identity identifies affine functions as the possible obstructions to density.  We now show that these are the only obstructions, using cutoffs in the $x$ coordinates.

\begin{lemma}\label{lem:density}
Let $d\mu=\rho(x)\,\rd x$ be a probability measure on $\mathbb{R}^n$ with a smooth, strictly positive, even density.  Let $\tau$ be a smooth even symmetric matrix field with $0<\tau(x)\le\Lambda I$ for some finite $\Lambda$, and set
\[
 Af=x\cdot\nabla f-\tau:D^2f.
\]
Suppose that $A$ has the integration formula
\[
 \mathbb{E}[hAf]=\mathbb{E}[\nabla h\cdot\tau\nabla f]
 \qquad(f,h\in C_c^\infty(\mathbb{R}^n))
\]
and satisfies the compact Hessian identity~\eqref{eq:hessian}. Then
\[
 \overline{\{A(A-1)f:f\in C_c^\infty(\mathbb{R}^n),\ f\text{ even}\}}^{\,L^2(\mu)}
 =\{g\in L^2(\mu):g\text{ even},\ \mathbb{E}g=0\}.
\]
\end{lemma}

\begin{proof}
Write $\Gamma(f,h)=\nabla f\cdot\tau\nabla h$ and $\Gamma(f)=\Gamma(f,f)$.  The integration formula also holds when one factor is smooth and the other has compact support, by inserting a cutoff around that support.  In particular $\mathbb{E}Ak=0$ for compactly supported smooth $k$.  Since $A$ preserves parity, the displayed range consists of centered even functions.

Let $g$ be a centered even $L^2(\mu)$ function orthogonal to this range. It also annihilates $A(A-1)f$ for odd $f$, because that image is odd. Formal symmetry therefore gives
\[
 \mathcal{P} g=0\quad\text{in distributions},\qquad \mathcal{P}=A(A-1).
\]
Indeed, if $\mathcal{P}^t$ denotes the transpose with respect to Lebesgue measure and $d\mu=\rho\,\rd x$, weighted formal symmetry gives $\mathcal{P}^t(\rho g)=\rho\mathcal{P} g$ in distributions. Orthogonality makes the left side zero, and smooth positivity of $\rho$ gives the displayed equation.  The principal symbol of $\mathcal{P}$ is $(\xi\cdot\tau(x)\xi)^2$, which is positive for every $\xi\ne0$. Interior elliptic regularity consequently gives $g\in C^\infty$ \cite[Theorem~14.2]{Dyatlov2026}.  Only positivity on compact sets is needed here.

\smallskip \noindent\emph{Splitting the equation in $L^2$.} Choose $0\le\eta_r\le1$ smooth, equal to one on $B_r$, zero outside $B_{2r}$, with $|\nabla\eta_r|\le C/r$ and $|D^2\eta_r|\le C/r^2$, where $r\ge1$.  The upper bound on $\tau$ gives constants $B,K$, independent of $r$, such that
\begin{equation}\label{eq:density-cutoff}
 \Gamma(\eta_r)^{1/2}\le B/r,
 \qquad |A\eta_r|\le K.
\end{equation}
For the second estimate, use $|x|\le2r$ on the support of the cutoff derivatives.  Those derivatives, including $A\eta_r$, vanish outside $B_{2r}\setminus B_r$.

Put $v=Ag$, initially just a smooth function, and write
\[
 G=\|g\|_2,\qquad
 X_r=\|\eta_r^2v\|_2,\qquad
 Y_r^2=\mathbb{E}[\eta_r^2\Gamma(g)].
\]
Here $G$ is finite by hypothesis, while $X_r,Y_r$ are finite by compact support. Testing $\mathcal{P} g=0$ against $\eta_r^4g$ gives $\mathbb{E}[v(A-1)(\eta_r^4g)]=0$. The product rule $A(ab)=aAb+bAa-2\Gamma(a,b)$ yields the exact expansion
\[
 \begin{split}
 (A-1)(\eta_r^4g)
 ={}&\eta_r^4v
 +g\bigl(4\eta_r^3A\eta_r-12\eta_r^2\Gamma(\eta_r)-\eta_r^4\bigr)\\
 &-8\eta_r^3\Gamma(\eta_r,g).
 \end{split}
\]
Thus Cauchy--Schwarz and~\eqref{eq:density-cutoff} imply
\begin{equation}\label{eq:density-two-estimates}
 X_r\le aG+\frac{8B}{r}Y_r,
 \qquad
 Y_r^2\le GX_r+\frac{2B}{r}GY_r,
 \qquad a=1+4K+12B^2.
\end{equation}
For the first inequality, factor out $\eta_r^2v$ in the testing identity. After this factorization, the term containing $g$ is bounded in $L^2$ by $aG$, and the last term is bounded by $(8B/r)\|\eta_r\Gamma(g)^{1/2}\|_2$. This uses $|\Gamma(\eta_r,g)|\le(B/r)\Gamma(g)^{1/2}$. Divide by $X_r$ if it is positive. For the second inequality, integrate by parts in $\mathbb E[\eta_r^2gAg]$:
\[
 Y_r^2=\mathbb{E}[\eta_r^2gv]-2\mathbb{E}[\eta_rg\Gamma(\eta_r,g)].
\]
Combining~\eqref{eq:density-two-estimates} gives $Y_r^2\le aG^2+10BGY_r$, so both $Y_r$ and $X_r$ are bounded uniformly in $r$.  Exhaustion by balls proves $v=Ag\in L^2(\mu)$.  We may now set
\[
 w=g-v\in L^2(\mu),\qquad Aw=0,\qquad Av=v.
\]

\smallskip \noindent\emph{The harmonic part is constant.} Testing $Aw=0$ with $\eta_r^2w$ and using \eqref{eq:density-cutoff} gives
\[
 \bigl(\mathbb{E}[\eta_r^2\Gamma(w)]\bigr)^{1/2}
 \le\frac{2B}{r}\|w\|_2.
\]
On any fixed ball, let $r\to\infty$.  Then $\Gamma(w)=0$ there; strict positivity of $\tau$ shows that $w$ is constant.

\smallskip \noindent\emph{The eigenvalue-one part is linear.} The analogous test of $Av=v$ gives, with $Z_r^2=\mathbb{E}[\eta_r^2\Gamma(v)]$,
\[
 Z_r^2\le\|v\|_2^2+\frac{2B}{r}\|v\|_2Z_r.
\]
Hence $\mathbb{E}\Gamma(v)<\infty$.  For $f_r=\eta_rv\in C_c^\infty$, the product rule now proves directly that
\begin{equation}\label{eq:density-approximation}
 f_r\longrightarrow v,\qquad Af_r\longrightarrow v
 \quad\text{in }L^2(\mu).
\end{equation}
Indeed,
\[
 Af_r=\eta_rv+vA\eta_r-2\Gamma(\eta_r,v).
\]
The middle term has norm at most $K\|v\mathbf{1}_{B_{2r}\setminus B_r}\|_2\to0$.  For the last term, use
\[
 \|\Gamma(\eta_r,v)\|_2
 \le\frac{B}{r}(\mathbb{E}\Gamma(v))^{1/2}\longrightarrow0.
\]
Apply~\eqref{eq:hessian} to these compactly supported functions:
\[
 \mathbb{E}\bigl[\|\tau^{1/2}D^2f_r\,\tau^{1/2}\|_{\mathrm{HS}}^2\bigr]
 =\|Af_r\|_2^2-\mathbb{E}[f_rAf_r]\longrightarrow0.
\]
On each fixed ball, $f_r=v$ for all sufficiently large $r$.  Positivity of the density and of $\tau$ therefore forces $D^2v=0$ on that ball. Thus $v$ is affine, and $Av=v$ removes its constant term.

We have proved that every $L^2$ distributional solution of $\mathcal{P} g=0$ is affine.  Our $g=w+v$ is even, so it is constant, and its mean is zero. Hence $g=0$.  Taking the orthogonal complement in the centered even subspace proves the lemma.
\end{proof}

% Adapted from an OpenAI Math manuscript. See PROVENANCE.md.
\section{A tensor estimate for the divergence tests}\label{sec:tensor}

The gradient and Hessian pairings in \eqref{eq:test-pairing} will be estimated using two quadratic energies, one for $W$ and one for $B$. The following lemma bounds their sum by $\mathbb{E}u^2$. Its proof uses weighted integration by parts and the moment equation; homogeneity enters only when the lemma is applied to the summands $\psi_i$ in \cref{sec:variance}.

\begin{lemma}\label{lem:tensor}
Let $d\mu=Z^{-1}\ee^{-V}dx$ and $d\nu=\ee^{-\varphi}dz$ be smooth probability measures on $\mathbb{R}^n$.  Suppose $D_z^2\varphi$ is positive definite and $x=\nabla\varphi(z)$ is a smooth diffeomorphism pushing $\nu$ to $\mu$. For a real $f\in C_c^\infty(\mathbb{R}_x^n)$, put
\[
 \tau=D_z^2\varphi,\quad M=\tau^{-1},\quad H=D_x^2V,
 \quad B=\tau(D_x^2f)\tau,
 \quad W=\delta_\mu B,
 \quad u=\delta_\mu W.
\]
View source fields in target coordinates by the inverse moment map. Then
\begin{equation}\label{eq:tensor-estimate}
 \mathbb{E}\tr\bigl((D_xW)^2\bigr)
 \ge \mathbb{E}\tr(MBHB),
\end{equation}
and
\begin{equation}\label{eq:tensor-energy}
 \mathbb{E}u^2\ge\mathbb{E}\tr(MBHB)+\mathbb{E}W^THW.
\end{equation}
Here $(D_xW)_{ij}=\partial_{x_j}W_i$, and the square in \eqref{eq:tensor-estimate} is a matrix square.
\end{lemma}

\begin{proof}
All test fields used in the integrations by parts below have compact support. Weighted Bochner calculations underlie many variance estimates; see Bakry--\'Emery \cite[Proposition~3]{BakryEmery1985} and the transport-Hessian formulas in \cite[Section~4.3]{KLR2026}. Here the required identity is for a vector field, so we derive it explicitly and retain the matrix trace $\tr((D_xW)^2)$. First, for any such vector field $W$ and $u=\delta_\mu W$,
\[
 \partial_{x_i}u
 =H_{ij}W_j+(V_{x_j}-\partial_{x_j})\partial_{x_i}W_j.
\]
Pairing with $W_i$ and using adjunction gives the weighted identity
\begin{equation}\label{eq:bochner}
 \mathbb{E}u^2
 =\mathbb{E}W_i\partial_{x_i}u
 =\mathbb{E}W^THW
   +\mathbb{E}(\partial_{x_j}W_i)(\partial_{x_i}W_j)
 =\mathbb{E}W^THW+\mathbb{E}\tr\bigl((D_xW)^2\bigr).
\end{equation}
Thus it remains to prove \eqref{eq:tensor-estimate}. Its proof compares two expansions, integrates their derivative terms by parts, and writes the resulting difference as a sum of squares.

Recall $h=D_x^2f$ and $P=\tau h$ from \eqref{eq:test-fields}, and let $(C_i)_{kl}=C_{ikl}$, where $C=D_z^3\varphi$. The moment equation \eqref{eq:moment} reads
\[
 V(x(z))=\varphi(z)+\log\det\tau(z)-\log Z.
\]
Its first derivative gives
\[
 \tau_{ik}V_{x_k}=x_i+c_i,
 \qquad c_i=\tr(MC_i).
\]
Since $\partial_jM=-MC_jM$, differentiating once more gives
\[
 (\tau H\tau)_{ij}
 =\tau_{ij}+M_{kl}\varphi_{klij}
   -\tr(MC_iMC_j)-C_{ijk}V_{x_k}.
\]
To combine the fourth-derivative term and the term involving $\nabla_xV$ into a weighted divergence, note that
\[
 \partial_dM_{dk}
 =-M_{da}C_{abd}M_{bk}=-c_bM_{bk},
 \qquad
 \partial_dM_{dk}-x_dM_{dk}=-V_{x_k}.
\]
Full symmetry of $D_z^4\varphi$ therefore yields
\begin{equation}\label{eq:twice-moment}
 (\tau H\tau-\tau)_{ij}
 =(\partial_d-x_d)(M_{dk}C_{ijk})
   -\tr(MC_iMC_j).
\end{equation}

For the other expansion, the commutator $[\partial_l,\partial_k^*]=\tau_{kl}$ gives
\[
 D_zW=B+R,\qquad R_{al}=\partial_k^*\partial_lP_{ak}.
\]
Both $D_zW$ and $B$ are symmetric, so $R$ is symmetric. Since $D_xW=(B+R)M$ and $MBM=h$, expansion followed by one integration by parts gives
\begin{equation}\label{eq:trace-expansion}
 \mathbb{E}\tr\bigl((D_xW)^2\bigr)
 =\mathbb{E}\tr(MBMB)
  +2\mathbb{E}(\partial_kh_{al})(\partial_lP_{ak})
  +\mathbb{E}\tr(MRMR).
\end{equation}
The last term is nonnegative: $\tr(MRMR)=\|M^{1/2}RM^{1/2}\|_{\mathrm{HS}}^2$. On the right side of \eqref{eq:tensor-estimate}, the coefficient of $\tau H\tau$ in the trace contraction is
\[
 Q=MBMBM=h\tau h.
\]
Substituting \eqref{eq:twice-moment} and integrating its divergence term by parts gives
\begin{equation}\label{eq:H-expansion}
 \mathbb{E}\tr(MBHB)
 =\mathbb{E}\tr(MBMB)
   -\mathbb{E}(\partial_dQ_{ij})M_{dk}C_{ijk}
   -\mathbb{E}Q_{ij}\tr(MC_iMC_j).
\end{equation}
The common first term cancels.  We now evaluate the three remaining contracted terms in coordinates where $\tau=I$ at the point under consideration.

This normalization uses a constant dual change of coordinates. For an invertible constant matrix $T$, set
\[
 z=Tz',\qquad x'=T^Tx,\qquad
 f'(x')=f(T^{-T}x').
\]
The corresponding normalized potentials and partition function are
\[
 \varphi'(z')=\varphi(Tz')-\log|\det T|,
 \qquad V'(x')=V(T^{-T}x'),
 \qquad Z'=|\det T|Z.
\]
They preserve both probability measures by change of variables, $x'=\nabla_{z'}\varphi'$, and the moment equation.  The chain rule gives
\[
 \begin{array}{lll}
 \tau'=T^T\tau T,& B'=T^TBT,& R'=T^TRT,\\
 M'=T^{-1}MT^{-T},& h'=T^{-1}hT^{-T},& H'=T^{-1}HT^{-T},\\
 Q'=T^{-1}QT^{-T},& P'=T^TPT^{-T},& W'=T^TW.
 \end{array}
\]
Each $z$-derivative index, including all three indices of $C$, transforms covariantly with $T$; the same law holds for $\partial_k^*$. Also $D_{x'}W'=T^T(D_xW)T^{-T}$. It follows that every contracted scalar in \eqref{eq:trace-expansion} and \eqref{eq:H-expansion} is invariant. For example, in $(\partial_kh_{al})(\partial_lP_{ak})$ each index pairs one covariant and one contravariant factor. At a fixed point we may therefore take $T=\tau^{-1/2}$, using the value of $\tau$ at that point.  All derivatives are computed under this constant transformation.  No moving frame is differentiated; the integrations by parts have already been completed.

In these normalized coordinates, put
\[
 S_{ijk}=f_{x_ix_jx_k},\qquad J_{ijk}=h_{ia}C_{ajk}.
\]
The tensor $S$ is fully symmetric, and $J$ is symmetric in its last two indices.  Write $J^{(12)}_{ijk}=J_{jik}$, and use the Euclidean tensor norm and inner product in this frame. The three contractions are
\begin{align}
 2(\partial_kh_{al})(\partial_lP_{ak})
   &=2\|S\|^2+2\langle S,J\rangle,\label{eq:normalized-cross}\\
 (\partial_dQ_{ij})C_{ijd}
   &=2\langle S,J\rangle+\langle J,J^{(12)}\rangle,
       \label{eq:normalized-derivative}\\
 Q_{ij}\tr(C_iC_j)&=\|J\|^2.\label{eq:normalized-cubic}
\end{align}
For \eqref{eq:normalized-cross}, use $\partial_kh_{al}=S_{alk}$ and $\partial_lP_{ak}=C_{ail}h_{ik}+S_{akl}$. For \eqref{eq:normalized-derivative}, differentiating the two factors $h$ in $Q=h\tau h$ gives the two copies of $\langle S,J\rangle$. Differentiating its middle factor gives
\[
 \sum_{iabjd}h_{ia}C_{abd}h_{bj}C_{ijd}
 =\sum_d\tr(hC_dhC_d)
 =\langle J,J^{(12)}\rangle.
\]
Finally, \eqref{eq:normalized-cubic} follows by expanding $Q=h^2$ and the trace, using the full symmetry of $C$.

Thus, apart from the nonnegative $R$ term, the integrand of \eqref{eq:trace-expansion} minus \eqref{eq:H-expansion} is
\[
 2\|S\|^2+4\langle S,J\rangle
   +\|J\|^2+\langle J,J^{(12)}\rangle.
\]
Full symmetrization of $J$ is
\[
 (\operatorname{Sym}J)_{ijk}
 =\frac{J_{ijk}+J_{jik}+J_{kij}}3.
\]
Its norm satisfies
\[
 \|\operatorname{Sym}J\|^2
 =\frac{\|J\|^2+2\langle J,J^{(12)}\rangle}{3},
 \qquad
 \langle S,\operatorname{Sym}J\rangle=\langle S,J\rangle.
\]
Consequently the preceding integrand equals
\begin{equation}\label{eq:tensor-squares}
 2\|S+\operatorname{Sym}J\|^2
 +\frac16\|J-J^{(12)}\|^2\ge0.
\end{equation}
This pointwise calculation proves \eqref{eq:tensor-estimate}; combining it with \eqref{eq:bochner} proves \eqref{eq:tensor-energy}.
\end{proof}

For the soft slab potentials, $H$ is the sum of a positive quadratic part and the positive semidefinite Hessians of the homogeneous summands.  \cref{lem:tensor} will control each summand through the two terms on the right side of \eqref{eq:tensor-energy}.

% Adapted from an OpenAI Math manuscript. See PROVENANCE.md.
\section{The variance inequality}
\label{sec:variance}

The tensor estimate controls two terms containing the Hessian of the potential. Each homogeneous summand contributes to both terms, and one Cauchy--Schwarz inequality combines its contributions with the correct coefficient. All expectations below refer to the probability $\mu$ of \eqref{eq:potential}.

\begin{proof}[Proof of \cref{prop:variance}]
The potential is smooth and even, and $D^2V\ge\varepsilon I$. \Cref{lem:moment} therefore supplies the moment coordinates, with $0<\tau\le (4/\varepsilon)I$. Put
\[
 H_i=D_x^2\psi_i,\qquad a_i=\mathbb{E}\psi_i.
\]
Each $H_i$ is positive semidefinite, and $a_i>0$ since $p_i\ne0$ and $\mu$ has positive density everywhere. Homogeneity gives
\[
 x\cdot\nabla\psi_i=q\psi_i,\qquad H_i x=(q-1)\nabla\psi_i.
\]
The moment-coordinate integration by parts gives
\[
 m_i:=\mathbb E\tr(\tau H_i)=\mathbb E(x\cdot\nabla\psi_i)=q a_i .
\]
For this integration by parts one can first multiply $\nabla\psi_i$ by a smooth cutoff equal to one on $B_R$ and supported on $B_{2R}$, with gradient $O(R^{-1})$. The cutoff error vanishes because $\tau$ is bounded and all polynomial moments are finite.

For even $f\in C_c^\infty(\mathbb R^n)$ put
\[
 B=\tau D_x^2 f\,\tau,\qquad W=\delta_\mu B,\qquad
 u=\delta_\mu W=A(A-1)f,\qquad M=\tau^{-1}.
\]
For each $i$, \eqref{eq:test-pairing} gives
\[
 d_i:=\mathbb E(u\psi_i)=\mathbb E(W\cdot\nabla\psi_i)=\mathbb E\tr(BH_i).
\]
In the $L^2(\mu)$ direct sum of matrices and vectors, with the Hilbert--Schmidt inner product for matrices, pair
\[
 (H_i^{1/2}\tau^{1/2},\,H_i^{1/2}x)
 \quad\text{with}\quad
 (H_i^{1/2}B\tau^{-1/2},\,H_i^{1/2}W).
\]
The pairing is $d_i+(q-1)d_i=q d_i$ by homogeneity. The squared norm of the first pair is $m_i+(q-1)m_i=q m_i$. Cauchy--Schwarz therefore gives
\[
 \frac{d_i^2}{a_i}
 \le \mathbb E\tr(MBH_iB)+\mathbb E(W^{\mathsf T}H_iW),
\]
using $m_i=q a_i$. Since $D^2V=\varepsilon I+\sum_iH_i$, \cref{lem:tensor} and \eqref{eq:bochner} bound the sum of the right-hand sides by $\mathbb E u^2$. It follows that
\begin{equation}\label{eq:dual-bound}
 \sum_i \frac{|\mathbb E(u\psi_i)|^2}{a_i}\le \|u\|_2^2 .
\end{equation}
The range density in \cref{lem:density} extends this inequality to every centered even $u\in L^2(\mu)$, since each $\psi_i$ belongs to $L^2(\mu)$.

For $F=\sum_i b_i\psi_i$ and centered even $u$ with $\|u\|_2\le1$, another Cauchy--Schwarz inequality gives
\[
 |\mathbb E(uF)|\le
 \left(\sum_i b_i^2a_i\right)^{1/2}
 \left(\sum_i|\mathbb E(u\psi_i)|^2/a_i\right)^{1/2}
 \le \left(\sum_i b_i^2a_i\right)^{1/2}.
\]
Taking the supremum in this subspace gives $\|F-\mathbb EF\|_2$ on the left. Squaring proves \eqref{eq:variance}, and \cref{sec:reduction} gives \cref{thm:main}.
\end{proof}

% Adapted from an OpenAI Math manuscript. See PROVENANCE.md.
\section{Volume and measure consequences}\label{sec:consequences}

With \cref{thm:main} proved, we derive the additive volume inequality stated in the introduction and then apply Saroglou's transfer theorem to even log-concave measures.

\subsection{The symmetric \texorpdfstring{$L_p$}{Lp} volume inequality}
\label{subsec:lp-proof}

\begin{proof}[Proof of \cref{cor:lp}]
Put $a=|K|^{1/n}$, $b=|L|^{1/n}$, $K_0=a^{-1}K$ and $L_0=b^{-1}L$. Set
\[
 c^p=(1-\lambda)a^p+\lambda b^p,qquad
 \theta=\lambda b^p/c^p.
\]
The weighted arithmetic--geometric mean inequality yields the pointwise bound
\[
 \bigl((1-\lambda)h_K^p+\lambda h_L^p\bigr)^{1/p}
 =c\bigl((1-\theta)h_{K_0}^p+\theta h_{L_0}^p\bigr)^{1/p}
 \geq c h_{K_0}^{1-\theta}h_{L_0}^\theta.
\]
The Wulff construction is monotone in its defining function and commutes with positive dilations. By \cref{thm:main}, its value on $h_{K_0}^{1-\theta}h_{L_0}^\theta$ has volume at least $1$. The left side of \eqref{eq:lp} is therefore at least $c^p$.
\end{proof}

\subsection{Even log-concave measures and the \texorpdfstring{$(B)$}{(B)}-conjecture}
\label{subsec:sec-measure}

We now pass from volume to integration against even log-concave measures. A nonnegative Radon measure $\mu$ on $\mathbb{R}^n$ is \emph{log-concave} if
\[
 \mu((1-\theta)A+\theta B)
 \ge \mu(A)^{1-\theta}\mu(B)^\theta
 \qquad(0<\theta<1)
\]
for all compact sets $A,B\subset\mathbb{R}^n$; it is \emph{even} if $\mu(A)=\mu(-A)$ for every Borel set $A$. Radon measures are finite on compact sets, but their total mass need not be finite. We use the analogous multiplicative definition of log-concavity for nonnegative functions, so the identically zero function is included. Endpoint products in an interpolation inequality are understood as the corresponding endpoint values.

Saroglou's transfer theorem gives the logarithmic Brunn--Minkowski inequality for an even log-concave density in any dimension in which it holds for Lebesgue measure \cite[Theorem~3.1]{Saroglou2016}. Averaging the measure over small translated balls will give such a density, even when the original measure is supported on a proper subspace.

\begin{corollary}
\label{cor:measure}
Let $n\ge1$, let $\mu$ be an even log-concave Radon measure on $\mathbb{R}^n$, and let $K,L\subset\mathbb{R}^n$ be origin-symmetric convex bodies. Then, for every $0\le\lambda\le1$,
\begin{equation}\label{eq:measure-logbm}
 \mu\bigl(\mathcal{W}[h_K^{1-\lambda}h_L^\lambda]\bigr)
 \ge \mu(K)^{1-\lambda}\mu(L)^\lambda.
\end{equation}
Consequently, for every origin-symmetric convex body $K\subset\mathbb{R}^n$, the function
\begin{equation}\label{eq:b-function}
 \mathbb{R}\ni t\longmapsto \mu(\ee^tK)
\end{equation}
is log-concave.
\end{corollary}

\begin{proof}
Fix $0<\lambda<1$, and put $D=\mathcal W[h_K^{1-\lambda}h_L^\lambda]$. Let $B$ be the closed unit ball. Choose $r>0$ with $rB\subset K\cap L$. The Wulff definition then gives $rB\subset D$ as well. We may restrict $\mu$ to a sufficiently large closed ball centered at $0$ containing $2K\cup2L\cup2D$. The restriction is finite by the Radon property. It is even and log-concave because the ball is symmetric and convex, and it has the same values on $K,L,D$. The zero measure and the endpoint cases give equality, so assume the restricted measure is nonzero.

Let $\mu_\varepsilon$ be its convolution with the uniform probability on $\varepsilon B$. Its density at $x$ is $f_\varepsilon(x)=\mu(x+\varepsilon B)/|\varepsilon B|$, by Fubini and symmetry of $B$. This density is even. It is log-concave by the defining inequality for $\mu$ on compact sets, since the Minkowski interpolation of $x+\varepsilon B$ and $y+\varepsilon B$ is the translate of $\varepsilon B$ by the interpolated center. Normalizing the density to have mass one leaves the transfer inequality unchanged, since both sides are homogeneous of degree one in the measure.

Take $0<\delta<1/2$ and $\varepsilon=\delta r$. If $x\in K$ and $z\in\varepsilon B$, then $x+z\in(1+\delta)K$ by convexity and $rB\subset K$. Thus $\mu_\varepsilon((1+\delta)K)\ge\mu(K)$, and the same holds for $L$. The Wulff interpolation of $(1+\delta)K$ and $(1+\delta)L$ is $(1+\delta)D$, by homogeneity. Apply \cref{thm:main} and Saroglou's transfer to get
\[
 \mu_\varepsilon((1+\delta)D)\ge
 \mu(K)^{1-\lambda}\mu(L)^\lambda.
\]
For a point $x+z$ counted by the left side, where $z\in\varepsilon B$, one has
\[
 x\in (1+\delta)D+\varepsilon B\subset(1+2\delta)D
\]
because $rB\subset D$. Therefore the left side is at most $\mu((1+2\delta)D)$. The dilates decrease to $D$ as $\delta\downarrow0$. Continuity of the finite measure from above gives \eqref{eq:measure-logbm}.

Finally $h_{\ee^t K}=\ee^t h_K$. The Wulff interpolation of $\ee^s K$ and $\ee^t K$ is therefore $\ee^{(1-\lambda)s+\lambda t}K$. Apply \eqref{eq:measure-logbm} to this pair to obtain \eqref{eq:b-function}.
\end{proof}


\printbibliography

\end{document}
